\section{The novelty filter in detail}\label{app:novelty}

This appendix gives Step~6 of \S\ref{sec:method} in full, and the sources searched for the outputs.

\subsection{The matching dataset of TheoremGraph}

The first pass is mechanical. Write a one-sentence slogan for the Lean statement of the rewritten theorem and look it up in the TheoremGraph matching dataset (dataset card, Hugging Face, accessed 23 November 2026) \cite{theoremgraph-matching}. Each of its 100,831 rows pairs a Lean declaration, from Mathlib or from one of 25 other projects, with the most similar arXiv statement, found by cosine similarity (at least 0.80) over embeddings of slogans; a language-model judge labels the pair \emph{exact}, \emph{inexact} or \emph{wrong}, and 47,952 pairs count as matches for its main judge. A hit shows that a formal twin exists and names the paper that states its informal version, and the statement goes to the calibration log with a precise reference. A miss shows nothing. An informal statement with no Lean declaration is not in the dataset, the labels are a model's and are not checked, and the informal text is included only for the pairs whose arXiv source allows redistribution, so the manual search is still needed. We ran this pass on eight outputs only (the quiver, cyclic-quiver, Morita, corner, $T_n$, $M_n(k)$, $(\mathbb Z/n,\cdot)$ and compact-Hausdorff statements; not the automata, the combinatorial topoi or the subspace lattices), by keyword search through the dataset's public search service, and no returned pair concerned any of them. We expected this, and it says nothing about novelty: the Lean side of the dataset has no declaration from Lean Pool and none about De Morgan topoi, and a keyword search ranks by words rather than by the similarity of slogans.

\subsection{The semantic search of TheoremSearch, and the OEIS}

For the informal side we ran the ten rows of the Outputs table through the semantic search of TheoremSearch \cite{theoremsearch}, which embeds a natural-language description and ranks some 8 million statements of arXiv and other sources by similarity of their slogans (public API, 23 November 2026, six results per query). No returned statement states any of the ten. Three queries returned ingredients that the corresponding sections already cite: the fundamental theorem of projective geometry, for the subspace lattices; for the compact Hausdorff output, the characterisation of De Morgan localic topoi (Johnstone's, in the form Caramello gives it) and the characterisation of compact Hausdorff spaces with Dedekind complete $C(X,\mathbb R)$ as the Stonean ones, as two separate statements; and, for corners, ring-theoretic statements in which a property of $R$ passes to $eRe$ when $ReR = R$, not the monoid statement. The other queries returned statements on quivers, cycle graphs, monoids or divisors with no bearing on the outputs. Six results per query make a shallow search, and a statement phrased in another vocabulary can be missed, so this pass narrows the manual search without replacing it. The calls, their parameters and the answers of 23 November 2026 are in \texttt{docs/bridges/novelty/} (\texttt{theoremsearch.py} and \texttt{theoremsearch-2026-11-23.json}): one request \texttt{POST https://api.theoremsearch.com/search} per output, with the description as \texttt{query}, \texttt{n\_results} set to 6, no filter, and no weighting by citations. The numerical outputs were also looked up in the OEIS \cite{oeis} (23 November 2026): the $2^n+1$ subtopoi of the $n$-cycle are A000051; the $2^{\omega(n)}$ essential points of $(\mathbb Z/n,\cdot)$-sets are A034444, the number of unitary divisors of $n$, as they should be, since the idempotents of $\mathbb Z/n$ correspond to its unitary divisors; the Dedekind numbers are A000372; the sequence $D(\omega(n))$ of subtopoi of $(\mathbb Z/n,\cdot)$-sets ($2, 3, 3, 3, 3, 6, 3, \dots$ for $n=1,\dots,16$) is not in the OEIS, although it is the composite of two sequences that are.

\subsection{The manual search}

The manual search covers, in this order: Johnstone's \emph{Elephant}; the papers of the author of the arch (here Caramello); the recent literature on the deck (here Hemelaer, Rogers, Pirashvili for monoid actions); the classical literature of the field of the rewritten statement (semigroups, quivers, topology); preprint servers and the nLab. Search for the \emph{rewritten} statement in the vocabulary of its field, not only for the topos-theoretic one: three of our candidates (the finite-monoid criterion, the subtopoi of $\mathbb N$-sets, essential points as classes of idempotents) were in print in the language of semigroups or monoid actions, and we first believed them new. Statements found in print go to the calibration log (\S\ref{sec:calibration}).

\subsection{Sources searched for the outputs}

For each statement of \S\ref{sec:machines}--\ref{sec:shapes} we searched Johnstone's \emph{Elephant}, Caramello's papers and book (including \cite{caramello-chamoun}), the papers and thesis of Hemelaer and Rogers on toposes of monoid actions \cite{hemelaer-rogers,rogers-1905,rogers-thesis}, the literature on Morita equivalence and Green's relations of monoids, the literature on synchronizing automata, the nLab, and general web and preprint searches. \q{Not found} means exactly that; several statements are elementary and may well be folklore, and we would welcome references.
